\documentclass[12pt,a4paper]{article}
\usepackage[utf8]{inputenc}
\usepackage[T1]{fontenc}
\usepackage{lmodern}
\usepackage[a4paper,margin=2.5cm]{geometry}
\usepackage{enumitem}
\usepackage{hyperref}
\usepackage{xcolor}
\usepackage{array}
\usepackage{booktabs}

\title{Projet Microservices\\Service A : Gestion des livres}
\author{Sirine}
\date{\today}

\begin{document}

\maketitle

\section{Introduction}

Dans le cadre du projet de microservices, ma partie porte sur le premier service de l'application : \texttt{books-service}.
Ce service a ete developpe en \textbf{JavaScript} avec \textbf{Node.js} et \textbf{Express}.

L'objectif de cette partie etait de realiser la progression suivante :

\begin{itemize}[leftmargin=*]
  \item developper un premier microservice en local ;
  \item conteneuriser ce service avec Docker ;
  \item publier son image sur Docker Hub ;
  \item deployer ce service sur Kubernetes ;
  \item ajouter une gateway locale avec Ingress.
\end{itemize}

\section{Presentation du service}

\texttt{books-service} est un microservice REST qui permet de gerer les livres d'une bibliotheque.

Fonctionnalites principales :

\begin{itemize}[leftmargin=*]
  \item ajouter un livre ;
  \item lister les livres ;
  \item consulter un livre par identifiant ;
  \item supprimer un livre.
\end{itemize}

\subsection{Routes exposees}

\begin{center}
\begin{tabular}{>{\ttfamily}p{3.5cm} p{9cm}}
\toprule
Route & Description \\
\midrule
GET /books & Lister tous les livres \\
GET /books/:id & Recuperer un livre par identifiant \\
POST /books & Ajouter un livre \\
DELETE /books/:id & Supprimer un livre \\
\bottomrule
\end{tabular}
\end{center}

\section{Etapes realisees}

\subsection{1. Service en local}

Le service a d'abord ete execute localement avec :

\begin{verbatim}
npm start
\end{verbatim}

La verification a ete realisee avec :

\begin{verbatim}
curl http://localhost:3001/books
\end{verbatim}

\subsection{2. Docker}

Un \texttt{Dockerfile} a ete cree pour conteneuriser \texttt{books-service}.

\begin{verbatim}
docker build -t books-service-local ./services/books-service
docker run --rm -p 3001:3001 books-service-local
\end{verbatim}

\subsection{3. Docker Hub}

L'image du service a ensuite ete publiee sur Docker Hub avec le compte :

\begin{verbatim}
sirinamarwa/books-service:latest
\end{verbatim}

Commandes utilisees :

\begin{verbatim}
docker tag books-service-local sirinamarwa/books-service:latest
docker push sirinamarwa/books-service:latest
\end{verbatim}

\subsection{4. Kubernetes}

Le service a ete deploye dans Minikube avec :

\begin{verbatim}
kubectl apply -f k8s/base/namespace.yaml
kubectl apply -f k8s/base/books-service.yaml
\end{verbatim}

Verification :

\begin{verbatim}
kubectl get deployments -n library-app
kubectl get pods -n library-app
kubectl get svc -n library-app
\end{verbatim}

Le pod \texttt{books-service} a bien ete observe dans l'etat \texttt{Running}.

\subsection{5. Gateway locale}

Une gateway locale a ete ajoutee avec un \texttt{Ingress} :

\begin{verbatim}
kubectl apply -f k8s/base/ingress.yaml
kubectl get ingress -n library-app
\end{verbatim}

La verification du routage a ete realisee avec :

\begin{verbatim}
kubectl port-forward -n ingress-nginx svc/ingress-nginx-controller 8081:80
curl -H "Host: library.local" http://localhost:8081/books
\end{verbatim}

\section{Etat actuel}

\begin{itemize}[leftmargin=*]
  \item service local valide ;
  \item Docker valide ;
  \item Docker Hub valide ;
  \item Kubernetes valide ;
  \item Ingress valide.
\end{itemize}

\section{Captures a inserer}

\begin{itemize}[leftmargin=*]
  \item lancement local du service ;
  \item reponse de \texttt{curl http://localhost:3001/books} ;
  \item construction de l'image Docker ;
  \item publication sur Docker Hub ;
  \item \texttt{kubectl cluster-info} ;
  \item \texttt{kubectl get nodes} ;
  \item \texttt{kubectl get pods -n library-app} ;
  \item \texttt{kubectl get svc -n library-app} ;
  \item \texttt{kubectl get ingress -n library-app} ;
  \item test Ingress via \texttt{curl -H "Host: library.local" http://localhost:8081/books}.
\end{itemize}

\section{Conclusion}

Cette premiere partie du projet valide toute la chaine technique pour un microservice :
developpement local, conteneurisation, publication d'image, deploiement Kubernetes et exposition via une gateway locale.

\end{document}
